\chapter{Longitudinal and survival software as a decision model}
\label{exec17:s5}

A bank comparing repeated servicing policies needs more than a curve with a
confidence ribbon. An account can default, repay early, cease to be observed,
or receive a different action after its financial condition changes. A policy
that reduces default incidence can also shorten the profitable relationship.
A policy that improves every baseline risk stratum can have a larger late
hazard among its remaining accounts. Two programs can return numerically
similar objects named survival while counting different accounts. These
differences affect both the scientific question and the financial decision.

The longitudinal identification and sequential-estimation results developed
earlier supply the foundation. Here the question is how to preserve that
foundation when choosing and combining software. The argument proceeds from
the event clock and supported horizon to survivor selection, competing-event
simulation, the meaning of a modified action, and uncertainty in the resulting
monetary contrast. Its finite laws and monetary amounts are stipulated,
reproducible constructions. They provide exact comparison references and
expose otherwise hidden distinctions; they are not bank observations or new
executions of the external estimators.\citep{exec17_s5_math}

\section{The decision clock and the unsupported tail}
\label{exec17:s5-horizon}

Let $j=1,\ldots,K$ index servicing periods. At the opening of a period the
bank observes the permitted history $H_j$, chooses an action $A_j$, and then
observes transitions and exits during that period. An intervention must specify
what the bank would do using the history available at that opening. A balance
posted after the decision, a later repayment classification, and a corrected
record unavailable to the decision maker cannot enter that history merely
because all three eventually appear in a warehouse. The data schema must
represent availability as well as the time to which a value refers.

Start with a single first-exit time $T$ taking values in the positive integers,
possibly with mass beyond the horizon. Under a specified policy $d$, define
$S^d(j)=P(T^d>j)$ and $S^d(0)=1$. Thus $S^d(j-1)$ is the probability of an
active account at the opening of period $j$. With payment for each period
entered, the expected restricted active time is
\begin{equation}
 R^d(K)=E\{\min(T^d,K)\}
       =\sum_{j=0}^{K-1}S^d(j).
 \label{exec17:s5-restricted-time}
\end{equation}
The identity follows account by account from
$\min(T,K)=\sum_{j=0}^{K-1}1\{T>j\}$ before taking expectations. This convention
earns a period's margin even if exit occurs at its end. A different payment
clock needs its own reward function. In particular, the same numerical area
must not be priced as both opening-period activity and completed-period
activity.

Suppose the observation and identification assumptions support the policy
survival probabilities only through period $m<K$. This may be an administrative
follow-up limit, rather than random censoring that an observation model could
remove. A fitted survival function beyond $m$ then supplies an extrapolation
assumption. Increasing simulation draws does not add observations at those
times. Without a tail restriction, monotonicity of survival gives the sharp
discrete bounds
\begin{equation}
 R^d(m)+S^d(m)\ \leq\ R^d(K)\ \leq\
 R^d(m)+(K-m)S^d(m).
 \label{exec17:s5-tail-bounds}
\end{equation}
To prove the lower bound, expand the remaining terms as
$S^d(m)+S^d(m+1)+\cdots+S^d(K-1)$ and use nonnegativity. To prove the upper
bound, bound every one of these $K-m$ terms by $S^d(m)$. The known term
$S^d(m)$ cannot be discarded. The continuous-time argument that allows all
remaining activity to disappear immediately after an observed time would
miss this whole discrete period.

Both endpoints are attainable. Keep the observed distribution through $m$
unchanged. For the lower endpoint, make all accounts still active at $m$ exit
at the end of period $m+1$. For the upper endpoint, keep them active through
the entire requested horizon. Mixing these two continuation laws preserves
every observed survival probability and fills the interval. The result is an
identified set conditional on the supported earlier survival law, not a
confidence interval for a point target. Sampling uncertainty about that earlier
law is a further problem.

For an exact six-period example, take the observed policy survival sequence
at times $0,1,2,3$ to be $(1,0.8,0.65,0.5)$ and the reference sequence to be
$(1,0.9,0.8,0.6)$. Then
\begin{align}
 R^d(6)&\in[2.95,3.95],& R^0(6)&\in[3.30,4.50],\notag\\
 R^d(6)-R^0(6)&\in[-1.55,0.65].
 \label{exec17:s5-tail-example}
\end{align}
These contrast endpoints are sharp when there is no additional cross-policy
restriction: select the lower continuation for one policy and the upper for
the other. Such potential event-time marginals can coexist, for example by
constructing each from a common uniform variable and its inverse distribution
function. A scientifically justified individual monotonicity restriction
would be additional information and could narrow the set.

At USD~20 margin per period entered and USD~3 upfront policy cost, the
incremental retention-margin value is therefore in $[-34,10]$ USD per eligible
account. This example prices activity alone; it does not silently include a
default loss or a repayment penalty. Those would require the corresponding
joint event probabilities. A decision to roll out the policy cannot be
justified by selecting the favorable end of this set. A bank with a more
restrictive accepted tail model could make a different, explicitly conditional
decision.\citep{exec17_s5_math}

There are several substantive remedies. Shortening the requested horizon to
the supported period changes the business question but removes this tail
ambiguity. Following a public or generated cohort longer can demonstrate how
a proposed analysis would use additional follow-up, without supplying the
bank's missing future. An externally justified survival envelope or a specified
parametric tail gives conditional bounds or a point extrapolation. The
report should show how the monetary conclusion changes across those choices.
An automatic default horizon chosen from the largest event in the current
sample answers a fourth question, whose end point varies with the sample.
That can be useful internally, but it does not establish comparability between
policies, software engines, or reporting vintages.

The software interface consequently needs an explicit requested horizon,
an observed-support horizon, an event convention and a monetary payment clock.
A response can legitimately contain a point estimate on the supported range
and a set beyond it. It must distinguish extrapolated uncertainty conditional
on a tail model from uncertainty that includes the choice of that model.
Returning a finite number at every requested time does not settle the
identification question.

\section{Why a late hazard can reverse a policy ranking}
\label{exec17:s5-selection}

Consider two equally weighted baseline risk strata, high and low, and no
competing exit or censoring. Within each stratum an active account defaults
with a constant probability per period. Under the reference policy these
probabilities are $0.8$ and $0.1$; under the intervention they are $0.4$ and
$0.05$. One can generate the potential histories using independent period
uniforms within each stratum. Randomized policy assignment, consistency and
the stated event clock then suffice for this illustrative policy comparison.
Its reversal does not require an unmeasured common cause of action and default.

Write $h_b^d$ for the hazard in baseline stratum $b$ and $\pi_b=1/2$ for its
baseline proportion. Marginal survival and the opening-period risk composition
are
\begin{align}
 S^d(j)&=\sum_b\pi_b(1-h_b^d)^j,\notag\\
 w_b^d(j-1)&=\frac{\pi_b(1-h_b^d)^{j-1}}{S^d(j-1)},\notag\\
 h^d(j)&=\sum_bw_b^d(j-1)h_b^d
        =1-\frac{S^d(j)}{S^d(j-1)}.
 \label{exec17:s5-mixture-hazard}
\end{align}
The two marginal hazards average over different survivor populations. By
period three, the reference high-risk share is $4/85$, whereas the policy
high-risk share is $144/505$. Protecting the high-risk accounts keeps more of
them in the later policy risk set. Comparing those risk sets as if they were
the same baseline population confuses composition with a policy's cumulative
consequence.

The exact results make the reversal visible:
\begin{center}
\begin{tabular}{lrr}
\hline
Quantity & Reference & Policy\\
\hline
$S(1)$ & $0.55$ & $0.775$\\
$S(2)$ & $0.425$ & $0.63125$\\
$S(3)$ & $0.3685$ & $0.5366875$\\
Period-three marginal hazard & $113/850$ & $1513/10100$\\
Three-period default risk & $0.6315$ & $0.4633125$\\
Restricted active periods & $1.975$ & $2.40625$\\
\hline
\end{tabular}
\end{center}
The period-three hazards are approximately $0.13294$ and $0.14980$. Thus the
policy has the larger late marginal hazard, although every stratum has a
smaller hazard under that policy at every time and marginal survival is larger
throughout. The marginal hazard ratio is a valid description of its two
selected risk sets. It is insufficient as a stand-alone ranking of the
baseline policy decision.

For a portfolio entering together, suppose each period entered earns USD~20,
default by period three loses USD~1,000, and the policy costs USD~3 upfront.
The incremental expected value is
\begin{equation}
 20(2.40625-1.975)
 -1000(0.4633125-0.6315)-3
 =173.8125\ \mathrm{USD}.
 \label{exec17:s5-selection-value}
\end{equation}
This ranking uses the same initial population and the cumulative event
probabilities appropriate to its financial objective. It remains conditional
on the stipulated margins and losses. Changing those amounts is a sensitivity
analysis of the objective, while changing how histories are selected is a
change to the population or causal model.\citep{exec17_s5_math}

This example motivates a practical output hierarchy. Preserve stratum-specific
and marginal hazards as diagnostics, but also report the policy survival or
cumulative incidence on a common baseline population, the restricted active
time and their joint uncertainty. If the bank wants an effect among accounts
that would survive under both policies, that is a principal-stratum question
requiring additional assumptions about joint potential histories. Neither
conditioning each fitted curve on its own observed survivors nor adjusting
the plot legend identifies that population.

The same logic affects model assessment. A regression can accurately predict
late defaults among observed reference survivors and still be a poor model
for the histories retained under another policy. Checking calibration on the
observed late risk set is informative about that set. Transporting it to a
policy-induced risk set requires overlap and stability of the relevant
conditional transition law. Flexible prediction and causal transport should
therefore appear as separate arguments in the proposal.

\section{Competing exits and the order of a simulator}
\label{exec17:s5-ordered}

Now distinguish default $D$ and early repayment $P$ as mutually exclusive
first events within each period. Among accounts event-free at the opening,
let $\lambda_{D,j}$ and $\lambda_{P,j}$ be their probabilities conditional on
the permitted history and intervention. They obey
$\lambda_{D,j}+\lambda_{P,j}\leq1$. With constant probabilities in a homogeneous
example,
\begin{align}
 S(j)&=(1-\lambda_D-\lambda_P)^j,\notag\\
 F_D(K)&=\sum_{j=1}^{K}S(j-1)\lambda_D,&
 F_P(K)&=\sum_{j=1}^{K}S(j-1)\lambda_P,\notag\\
 S(K)+F_D(K)+F_P(K)&=1.
 \label{exec17:s5-first-events}
\end{align}
Here $S$ means free of both first events. In contrast,
$1-F_D=S+F_P$ includes accounts that have already repaid. It cannot be
integrated as if those repaid accounts were still earning the active-account
margin. Treating prepayment as censoring would target another quantity and
needs a separate scientific interpretation; it does not automatically
describe eliminating prepayment while keeping every other mechanism fixed.

Some simulators implement within-period draws sequentially. If prepayment is
drawn first with probability $p_P$ and default is then drawn among those
without that prepayment with probability $p_D$, the induced unconditional
first-event probabilities are
\begin{equation}
 \lambda_P=p_P,\qquad \lambda_D=(1-p_P)p_D,
 \qquad P(\text{neither})=(1-p_P)(1-p_D).
 \label{exec17:s5-ordered-conversion}
\end{equation}
Consequently an adapter from mutually exclusive first-event probabilities
must set $p_P=\lambda_P$ and
$p_D=\lambda_D/(1-\lambda_P)$ when $\lambda_P<1$. The constraint on the
two first-event probabilities ensures $p_D\leq1$. If $\lambda_P=1$, feasibility
requires $\lambda_D=0$ and the unused conditional default probability does not
determine any outcome. Reversing the order requires the converse
factorization. Multiplying two arbitrary competing hazards without defining
their conditioning populations supplies neither factorization.

The conversion is elementary but consequential. Consider two periods. Set
$(\lambda_D,\lambda_P)=(0.1,0.1)$ for reference and $(0.12,0.6)$ for policy.
Although the policy's default probability among initially event-free accounts
is larger each period, more of its accounts prepay before a later default is
possible. The resulting $(F_D,F_P,S)$ vectors at period two are respectively
$(0.18,0.18,0.64)$ and $(0.1536,0.768,0.0784)$. The policy has the smaller
default cumulative incidence and the shorter active relationship.

Price each entered period at USD~20, default at a USD~1,000 loss, prepayment
at a USD~50 loss, and the intervention at USD~3 per initial account. The
joint objective is
\begin{equation}
 V^d(K)=20R^d(K)-1000F_D^d(K)-50F_P^d(K)-c_d.
 \label{exec17:s5-event-value}
\end{equation}
Using $c_0=0$ and $c_d=3$ gives reference value $-153$ USD and policy value
$-169.4$ USD. The incremental value is $-16.4$ USD despite a smaller default
incidence. Default avoidance alone does not establish a desirable servicing
policy when relationship duration and competing exits carry monetary value.

For the ordered simulator the correct default probabilities after a
non-prepayment draw are $1/9$ for reference and $0.3$ for policy. Suppose an
adapter instead passes the policy's unconditional $0.12$ directly as $p_D$,
while setting $p_P=0.6$. The induced period default probability is then
$0.4\times0.12=0.048$. At period two the wrong policy calculation returns
\begin{align}
 F_D^{\rm wrong}&=0.064896,&F_P^{\rm wrong}&=0.8112,\notag\\
 R^{\rm wrong}&=1.352,&V^{\rm wrong}&=-81.416\ \mathrm{USD}.
 \label{exec17:s5-ordered-error}
\end{align}
Compared with the correctly translated reference, the apparent gain becomes
$71.584$ USD. An omitted conditional-probability conversion has reversed the
decision. The example does not identify a defect in an upstream package: the
error is the deliberately incompatible meaning supplied at its boundary.
\citep{exec17_s5_math}

The inspected simulation path of pygformula draws a competing event before
the target event and marks the target event missing after that competing
exit. Its aggregation multiplies the target-event probability by the
no-competing-event probability and the preceding survival products. The
stored object named survival is one minus accumulated target incidence;
it is therefore distinct from event-free survival when competing exits are
present. These are claims about the inspected original simulation and
aggregation paths, not about newly fitted models or all package releases.
The model-fitting risk sets were not newly inspected for this comparison,
so the simulation order alone cannot establish how every fitted probability
is conditioned.\citep{v13_sw_pygformula,exec17_s5_paths}

An integration should first check the complete tuple: opening risk set,
within-period order, event codes, probability conditioning, transition to an
absorbing state, and the meaning of every accumulated output. A generated
two-period law is a useful exact semantic test before statistical fitting.
It can verify that a proposed wrapper conserves probability and obtains
Equation~\eqref{exec17:s5-event-value} under declared inputs. It does not
establish that the observational data identify those inputs, or that a
fitted observational model is calibrated under intervention.

\section{Modifying a natural action and redrawing an action}
\label{exec17:s5-natural}

A dynamic servicing policy may modify the action that would naturally be
chosen: shorten a delay, reduce a proposed payment, or add support when the
proposed contact intensity lies in a specified range. This is more specific
than drawing a new action from a distribution with the same conditional
marginal frequencies. At later times the natural action is evaluated under
the history produced by earlier interventions. Unobserved features of the
action mechanism can connect that later natural action to an earlier one,
even when there is no unmeasured confounding of the ordinary joint treatment
regime and outcome.

The longitudinal modified treatment policy analysis distinguishes a natural
action policy from its stochastic-intervention counterpart. Under its
structural model, standard sequential randomization separates a current
treatment innovation from later covariate and outcome innovations. The
stronger condition also separates it from future treatment innovations.
Theorem~1 identifies the redrawn stochastic intervention under the standard
condition and the natural modified policy under the stronger condition,
together with support. The two targets must therefore be named before
reusing the same observed-data recursion.\citep{lmtp2023}

An exact two-time construction shows why this distinction is necessary.
Let $U$ and $V$ be independent fair bits, with no measured covariates. In both
worlds $A_1=U$ and the structural outcome is $Y=A_2$. The second natural action
is generated by
\begin{equation}
 \begin{array}{ll}
 \mathcal W_0:& A_2=V,\\
 \mathcal W_1:&
 A_2=\begin{cases}V,&A_1=U,\\1,&A_1\ne U.\end{cases}
 \end{array}
 \label{exec17:s5-two-worlds}
\end{equation}
In ordinary observations $A_1=U$, so the second branch of $\mathcal W_1$ is
never used. In each world all four $(A_1,A_2)$ cells have probability $1/4$,
and $Y=A_2$. Thus the complete observable law is identical, with positive
support for both actions at each observed history. This is equality of the
whole joint distribution, not merely equality of observed outcome means.

Apply the policy $d_1(a)=1-a$, followed by $d_2(a)=a$: flip the first natural
action and leave the second natural action as generated under the changed
history. In $\mathcal W_0$, the second action remains $V$ and
$E(Y^{\bar d})=1/2$. In $\mathcal W_1$, the first action becomes $1-U$;
the second action mechanism then takes its other branch and sets $A_2=1$.
Its modified-policy mean is $1$. The natural-action effect is not identified
by the common observed law under the weaker restrictions of this construction.

For the redrawn intervention, the observed conditional action law at both
times is a fair bit. Pushing those conditional laws through the flip and
identity maps still gives fair bits; drawing from them under the evolving
history yields outcome mean $1/2$ in both worlds. This is a different
intervention. It is identified by the common observed law under the stated
standard condition. Calling it a numerical approximation to the natural
policy would obscure a difference in counterfactual meaning.

The construction satisfies the standard innovation condition used above:
covariate and outcome structural innovations can be constants, and the
outcome under a joint assigned sequence is simply the assigned second
action. The second treatment innovation in $\mathcal W_1$, however, contains
the earlier $U$. Strong independence between the first and later treatment
innovations fails. In $\mathcal W_0$ independent $U,V$ suffice. Support also
holds in the example; the failure does not arise from a missing observed
action cell. These details explain why the witness illustrates the theorem's
assumption distinction rather than contradicting its identification result.
\citep{lmtp2023,exec17_s5_math}

In banking terms, a latent servicing preference that influences several
successive proposed actions can matter to a policy that modifies those
proposals. Observing all outcome confounders for ordinary assigned regimes
does not automatically reveal that proposal mechanism. This interpretation
is a possible mechanism, not an assertion about a bank's data. A proposal
should ask whether the natural recommendation can be recorded, whether its
history-dependent generation can be justified, and whether a randomized
redraw is operationally acceptable. If the business instead intends a fresh
randomized allocation at every time, the stochastic target may be exactly
the intended one.

There are consequently several defensible responses to uncertain natural
action dependence: record and experimentally vary the recommendation
mechanism; impose and defend the stronger structural restriction; redesign
the intervention as a redrawn policy; or retain an explicitly incomplete
identification claim. An estimator's robustness properties concern an
identified observed-data functional under its assumptions. They cannot
choose among observationally identical structural worlds with different
natural-policy outcomes.

\section{A supported shift, a cap, and the boundary of the density ratio}
\label{exec17:s5-pushforward}

Continuous treatment policies create a second boundary that is easy to hide in
an API. Let $A\sim\operatorname{Uniform}[0,1]$ and consider the fixed map
\begin{equation}
 d(a)=\begin{cases}a+1/4,&0\leq a\leq3/4,\\a,&3/4<a\leq1.\end{cases}
 \label{exec17:s5-shift}
\end{equation}
The map is not globally one-to-one. Its image has two inverse branches on
$[3/4,1]$: the identity branch and the shifted branch. The pushforward density
relative to Lebesgue measure is therefore
\begin{equation}
 g_d(y)=\begin{cases}
 0,&0\leq y<1/4,\\
 1,&1/4\leq y<3/4,\\
 2,&3/4\leq y\leq1.
 \end{cases}
 \label{exec17:s5-push-density}
\end{equation}
The value two is not a probability exceeding one. It is the sum of the masses
from two preimages. Dropping one inverse branch changes the intervention law
and its likelihood ratio. Integrating the density gives one, and direct
integration gives
\begin{equation}
 E\{d(A)\}=11/16,\qquad E\{d(A)^2\}=25/48,qquad
 \operatorname{Var}\{d(A)\}=37/768.
 \label{exec17:s5-push-moments}
\end{equation}

Now compare the apparently harmless cap $d_c(a)=\min(a+1/4,1)$. It maps the
interval $(3/4,1]$ to the single point one, creating an atom of mass $1/4$.
The observed continuous treatment law gives that point probability zero, so
the intervention law is not absolutely continuous with respect to the observed
law. A density-ratio estimator defined only with respect to the observed
Lebesgue density cannot represent this target. This is a support failure, not
merely a large finite weight.

The distinction matters even if an outcome regression is available. Suppose
$m(a)=E(Y\mid A=a)$ is known only as a version up to Lebesgue-null sets and
the outcome is bounded in $[0,1]$. The value at the cap point is then
unrestricted by the observed regression. The target mean lies in
\begin{equation}
 \left[\int_0^{3/4}m(a+1/4)\,da,
       \int_0^{3/4}m(a+1/4)\,da+1/4\right]
 = [15/32,23/32]
 \label{exec17:s5-cap-set}
 \end{equation}
for the exact authored regression $m(a)=a$ on the observed continuum. If the
outcome is worth USD~100 at its unit value, the baseline is USD~50 and the
intervention costs USD~3, the net value set is
$[-49/8,151/8]$ USD. A continuous extension of $m$ to the endpoint, a structural
model for the atom, or a redesigned policy that stays in the observed support
would add information. Simply clipping the action and retaining the continuous
density ratio would silently select one endpoint or an incompatible target.
\citep{lmtp2023,exec17_s5_math}

For the supported fixed map in Equation~\eqref{exec17:s5-shift}, write
$g(a\mid W)$ for the observed conditional treatment density,
$r(a,W)=g_d(a\mid W)/g(a\mid W)$, and $m(a,W)$ for the conditional outcome
regression. The one-time target is
\begin{equation}
 \theta(P)=E_W\int m(d(a,W),W)g(a\mid W)\,da.
 \label{exec17:s5-lmtp-target}
\end{equation}
The fixed-map path derivative has the influence function
\begin{equation}
 D(P)(O)=r(A,W)\{Y-m(A,W)\}+m(d(A,W),W)-\theta(P).
 \label{exec17:s5-lmtp-if}
\end{equation}
The map is fixed independently of the law $P$, although its induced density
$g_d$ changes when $g$ changes. If the assignment law is instead frozen as an
external design, the assignment contribution is averaged under that fixed law:
\begin{equation}
 M(W)=\int m(d(a,W),W)g(a\mid W)\,da,qquad
 D_{\mathrm{frozen-q}}=r(A,W)\{Y-m(A,W)\}+M(W)-\theta.
 \label{exec17:s5-frozen-if}
\end{equation}
The extra random term in the observational fixed-map score is
$m(d(A,W),W)-M(W)$. It is a design component, not a nuisance-residual
component. Confusing the two scores changes the variance claim.

The exact check uses the shift in Equation~\eqref{exec17:s5-shift} and
$m(a)=a$. Along the normalized assignment path
$g_\epsilon(a)=1+\epsilon(2a-1)$, the derivative at zero is $23/192$ and
$\theta(1/4)=551/768$. For a Bernoulli outcome with conditional mean $A$, the
residual variance is $7/32$, the assignment component is $37/768$, and the
full score variance is $205/768$. These values are algebraic checks of the
target and score under a declared law, not an empirical efficiency result.
The fixed-map and piecewise-invertibility conditions in the primary source are
therefore operational requirements for an adapter. The source's threshold
example also illustrates that a point intervention at a continuous boundary
can be non-pathwise-differentiable. A software response should expose this
boundary instead of reporting a routine root-$n$ interval.\citep{lmtp2023,exec17_s5_math}

\section{Translating software semantics before comparing estimates}
\label{exec17:s5-software}

The common package names in a longitudinal workflow conceal different data
contracts. The inspected ltmle documentation uses wide, time-ordered nodes;
an absorbing event is encoded as event one, censoring is a separate factor,
and post-event or post-censor nodes are ignored. Its documented bounds act on
cumulative treatment probabilities. A continuous outcome range can also be
truncated or scaled. These choices determine the observed target before a
targeting step is run.\citep{sw_ltmle}

The pinned source adds an operational uncertainty boundary. A requested
`variance.method` other than `ic` is reset to `ic` when a non-null cluster
identifier is supplied; the effect-summary path then uses the influence-curve
covariance, while the inspected household aggregation first sums record-level
contributions within each identifier. Thus the identifier changes the sampling
unit of that calculation, but does not preserve the package's additional
TMLE-based variance option. A request can therefore return without the
variance method its caller named. Any adapter must expose the effective method,
the realized cluster count and cluster-size distribution, and the unit used to
scale the covariance. The source permits only binary treatment nodes, and the
single-regime or two-regime `ltmle` wrapper is distinct from an MSM fit for a
larger set of policies. These are interface facts, not claims that its
influence-curve estimator is or is not calibrated for a particular bank
sampling law.\citep{sw_ltmle}

For example, if each household contributes two identical account records with
influence value $\psi$, the household contribution is the sum $2\psi$ before
the package rescales it from records to households. An iid-row calculation
would divide by twice as many nominal units and understate the variance of the
household-mean estimator by a factor of two. Correct aggregation addresses
that sampling-unit error; it does not establish independent households,
correct censoring models, adequate positivity or calibrated finite-sample
coverage.

An exact authored two-period example makes the policy target concrete without
pretending to run the package. Let $H$ mark a high-risk account,
$P(H=1)=1/3$, and let $L_1$ be subsequent hardship. Under the policy's first
action $A_0$, set
$P(L_1=1\mid H,A_0)=(1/4,1/8,3/4,1/2)$ for
$(H,A_0)=(0,0),(0,1),(1,0),(1,1)$ respectively. Conditional default risk is
$1/20+H/10+L_1/5-A_0/25-A_1/20$. Compare status quo
$(A_0,A_1)=(0,0)$, the dynamic targeted rule $A_0=H,A_1=L_1$, and broad
contact $(A_0,A_1)=(1,1)$. Censoring leaves an account observed with probability
$9/10$ for $L_1=0$ and $3/4$ for $L_1=1$. Exact path enumeration gives
\begin{align}
 \theta_d &= \sum_{h,l}P(h)P(l\mid h,A_0=d_0(h))
   E[Y\mid h,A_0=d_0(h),l,A_1=d_1(l)],\\
 R(d) &=1000\{\theta_{\mathrm{status}}-\theta_d\}
       -6E_d(A_0)-4E_d(A_1).
\end{align}
The default probabilities for status quo, targeted and broad contact are
$1/6$, $3/25$ and $13/300$; their net changes in expected value against
status quo are respectively USD $0$, $130/3$ and $340/3$ per account. All
these values are exact under the stipulated law. The minimum observation
probability is $3/4$, but that is not a sampling guarantee beyond this law.
The calculation does not estimate censoring weights or uncertainty, execute
ltmle, or establish that broad contact is beneficial in bank data. It tests
whether a proposed workflow preserves the dynamic rule, regime mean, censoring
history and the separate decision-layer cost.

The lmtp documentation uses chronological treatment, covariate and outcome
nodes. A supplied shifted rule overrides a generic shift, finite memory limits
which history enters the intervention, and cluster identifiers, folds and
weights affect the fitted calculation. Its returned density ratios are
noncumulative. A wrapper that turns them into cumulative ratios, or clips a
cumulative ratio at each time instead of clipping the final product, has
changed both numerical uncertainty and often the estimand.\citep{sw_lmtp}

For a two-step exact probability law with incremental denominators
$(1/2,2/5)$, the incremental ratios are $(2,5/2)$ and their cumulative
products are $(2,5)$. If a floor of $1/4$ is applied to the cumulative
denominator, the final weight is $4$; if the floor is applied to each
incremental denominator before multiplying, the final weight is $5$. Both
procedures avoid an infinite value, but neither is automatically target
preserving. The floor level, location, and resulting target must be recorded
as part of the policy definition.

The stremr manual presents a long-data interface with subject and time fields
and separates treatment, monitoring and censoring models. Its `fitTMLE` path
wraps sequential g-computation with targeting. The survtmle manual fixes a
binary treatment, integer event times and an event code for censoring; its
default horizon is derived from the observed event time. It documents mean and
hazard methods, covariance and subject-level influence vectors, and warns that
finite-sample convergence is not guaranteed. The prose and argument material
also use different descriptions of a tolerance scaling; an adapter must pin
the actual installed behavior rather than infer it from prose.\citep{v12_sw_stremr,v12_sw_survtmle}

The gfoRmula API exposes risk and end-of-follow-up outcomes, censoring and
competing-event options, and custom intervention and history roles. Its names
do not by themselves establish whether an output is event-free survival,
cumulative incidence, or a different complement. The pygformula simulation
path draws competing events before target events, while its bootstrap resamples
whole histories for unique subjects, refits models, and uses common seeds for
intervention simulations. Those are useful building blocks for a prospective
comparison. They do not prove that every model-fitting risk set or every
release has the same semantics.\citep{sw_gformula,v13_sw_pygformula}

An interoperable result should therefore carry at least: subject unit and
cluster, clock and horizon, event and censoring codes, absorbing-state rule,
policy law and whether it is natural or redrawn, treatment probability scale,
ratio orientation and truncation, target type, influence-vector convention,
failed-replicate handling, software version and license. A result lacking one
of these fields can remain a useful exploratory number, but it cannot enter a
common-target comparison without an explicit refusal or translation proof.
Package declarations currently record gfoRmula 1.1.2 (GPL-3), ltmle 1.3-0
(GPL-2) and lmtp 1.5.4 (AGPL-3.0). The inspected API bytes and package
declarations are separate evidence; metadata alone does not prove that those
bytes belong to those exact releases. Complete transitive dependency rights,
runtime qualification and bank acceptance remain open.

\section{Sampling, simulation and model error are different quantities}
\label{exec17:s5-uncertainty}

Suppose a fitted model generates a policy value by drawing a baseline history
$D$, then Monte Carlo transitions $U$. For an estimator based on $M$ simulated
histories, the law of total variance gives
\begin{equation}
 \operatorname{Var}(\widehat V)=
 \operatorname{Var}_D\{E(\widehat V\mid D)\}+
 E_D\{\operatorname{Var}(\widehat V\mid D)\}.
 \label{exec17:s5-total-variance}
\end{equation}
The first term is variation in the fitted baseline and observed sample; the
second is simulation noise conditional on that fitted object. Model
misspecification is a bias term and is not removed by increasing $M$. A
reported Monte Carlo standard error that omits the first term is not a
confidence interval for the full causal estimand.

For a paired policy contrast, use one common uniform draw per simulated unit
and set the event indicator to $1\{U<p_d\}$. With $p_0=0.2$ and $p_1=0.3$, the
single-draw difference has mean $0.1$ and variance $0.09$, because it is one
when $U\in(0.2,0.3]$ and zero otherwise. Independent draws have variance
$0.2(0.8)+0.3(0.7)=0.37$. At $M=10{,}000$, the corresponding simulation
variances are $0.000009$ and $0.000037$. Common random numbers can improve a
contrast by coupling comparable paths, but resetting a seed alone does not
guarantee that branching histories remain aligned. The coupling rule must be
declared.

If each independent account contributes two identical rows, treating rows as
iid understates the variance of a subject-level mean by a factor of two. The
bootstrap and influence calculation must resample or retain the independent
subject unit. The pygformula bootstrap's whole-history resampling is therefore
semantically different from a row bootstrap. A future package comparison must
record this unit explicitly.\citep{v13_sw_pygformula,exec17_s5_math}

Coverage also needs a complete denominator. In an illustrative 200-request
study, 160 requests are available, 152 are covered, so availability is $0.8$,
conditional coverage is $152/160=0.95$, and request coverage is $152/200=0.76$.
Reporting only the conditional fraction would hide the failed or unavailable
requests. Pointwise intervals for several horizons or policies are not
simultaneous bands; retaining the joint subject influence vectors permits a
declared covariance calculation, while simultaneous control needs an explicit
multiplier or resampling rule. These are reporting choices with direct business
meaning when a decision is based on the worst horizon or the best policy.

\section{A bounded transition-error envelope}
\label{exec17:s5-error}

An integration can state a numerical error budget without pretending to know a
bank model's predictive accuracy. Let $P_t(\cdot\mid h)$ and
$\widetilde P_t(\cdot\mid h)$ be the true and implemented transition kernels
for histories at step $t$, and suppose their total-variation distance is at
most $\varepsilon_t$ uniformly over admissible histories. Couple the two
processes maximally at each step. If $d_t$ bounds the probability that the
histories have diverged before step $t$, then kernel contraction and a union
bound give
\begin{equation}
 d_0=0,\qquad d_{t+1}\leq\min\{1,d_t+\varepsilon_t\}.
 \label{exec17:s5-tv-recursion}
\end{equation}
For a running reward with range at most $B_t$ at step $t$, the absolute
expected running-reward bias is at most
\begin{equation}
 \sum_{t=0}^{K-1}B_t d_t,
 \label{exec17:s5-running-bound}
\end{equation}
and a terminal loss with range at most $L$ adds at most $Ld_K$. To see this,
condition on the common history event and use equality of coupled rewards
there; on the divergence event the reward difference is bounded by its range.
This is a proof under the declared uniform TV budget, not a measured error
rate or a predictive MSE guarantee.

For $K=3$, $arepsilon_t=0.02$, $B_t=20$ USD and terminal loss range
$L=1000$ USD, the recursion gives
$d=(0,0.02,0.04,0.06)$. The running-margin bound is USD~1.20 and the terminal
default bound USD~60, for a total of USD~61.20. A bank can replace these
uniform budgets with validated conservative budgets only after obtaining
appropriate transition audits. A small average error on frequent histories
does not imply this uniform bound on rare but financially important histories.
Likewise, a terminal value range should include all competing event outcomes,
not only defaults, if the policy objective prices them jointly.

The envelope offers a useful admission test for an implementation. The
adapter must expose the horizon, transition decomposition, reward range,
stated TV or coupling budget, and whether the bound includes model fitting,
simulation, and observation error. If any budget is unavailable, the system
can return a conditional point estimate with a declared numerical component,
but it should not label the resulting interval a full causal uncertainty
interval. Refusal is a valid output when the requested guarantee exceeds the
declared budget.

\section{A common-target public prototype and adoption routes}
\label{exec17:s5-adoption}

The practical next step is a public or generated comparison with a common
target, not a claim that six packages are interchangeable. Freeze a finite
baseline cohort law, an integer event clock, explicit default/repayment/censor
codes, two policies, a supported horizon, monetary rewards, and a complete
seed and environment manifest. Admit an engine only after exact two-period and
single-shift references reproduce the probability, pushforward and reward
checks above. A candidate that cannot express an atom, competing first event,
natural-versus-redrawn choice, or subject-level resampling is not silently
coerced into the target; it receives a typed incompatibility record.

For each admitted engine, register the input schema, version and license,
training folds, policy specification, simulation coupling, number of baseline
and Monte Carlo draws, interval or band method, failed-replicate rule and
runtime. Compare the policy contrast, restricted active time, each first-event
incidence, joint covariance, availability, width, numerical error and the
resulting monetary decision. Keep model bias as a separate unknown. A generated
potential-outcome law establishes the truth of that construction only. Public
data can exercise the observation and clock pipeline, but cannot establish a
bank's unmeasured assignment truth or accepted loss threshold.

Adoption then has four defensible routes. Direct adoption is suitable when an
engine's target, event semantics, license and runtime are accepted. A thin
adapter is preferable when the estimator is sound but its schema or result
metadata differ. Porting is justified only when the missing capability is
material, the license permits it, and a reference implementation can be
reproduced. Rejection is appropriate for an unsupported target, unresolvable
license or dependency risk, missing uncertainty meaning, or an output that
cannot preserve the bank's information clock. These are proposal decisions,
not claims that the underlying packages are defective.

The shared library should expose a typed analysis manifest and a refusal
object alongside estimates. The CLI can validate the manifest, run exact
semantic fixtures, and emit a versioned result bundle. An MCP surface can
request validation, explanation, and comparison while preserving the same
manifest and audit record; it should not infer a causal diagram or silently
choose a horizon from an unqualified natural-language request. LLM or agent
assistance can propose mappings and diagnostics, but a human or registered
rule must approve the target, assumptions, source clock and policy. The
longitudinal engine should return an identified set or an incompatibility when
the tail, atom, event order or natural-policy condition is unsupported.

This design is ready for a prospective public/generated pilot and for package
wrappers under the current proposal authorization. It does not close the
software census, complete dependency-rights review, independent scientific
acceptance, authentic bank validation, target-platform recovery, or economic
adoption. Those gaps require external inputs and remain explicit in the master
program.\citep{exec17_s5_design}

\section{Conclusion}
\label{exec17:s5-end}

Longitudinal causal software is comparable only after the decision clock,
event law, intervention meaning, support, uncertainty unit and monetary
objective are aligned. The exact constructions show where a plausible scalar
can reverse a decision: an unsupported tail, survivor composition, an ordered
competing-event conversion, a natural action that is not a redraw, a singular
cap, a cumulative-weight translation, or a simulation error presented as a
confidence interval. A product proposal can close these semantic and design
gaps with public/generated fixtures and bounded derivations now. It cannot
claim that those fixtures establish bank truth, general calibration, or final
adoption.
